\section{问题重述}

\subsection{问题背景}
微网是消纳光伏等分布式能源、减少并网冲击的重要形式。由于光伏出力具有明显的间歇性与波动性，
小区负载也随时间变化，非孤岛式微网可在外部电网（下称外网）电价较低、或分布式能源有剩余电量的
时段为储能设备充电，并在电价较高的时段释放电能，从而在保证供电可靠的同时降低购电成本。
题目给出了某微网一天的分时电价、小区负载与光伏发电预测功率，以及储能设备的运行参数：
最大容量 12000\,kWh、最大充放电功率 5000\,kW、充放电效率 90\%，
且运行过程中储电量须保持在 1200$\sim$10800\,kWh 之间。
在此基础上，需要针对电价、负载与光伏信息在不同时间尺度上的已知程度，
分别建立微网日前购电与储能充放电的优化调控模型，并给出相应的计算结果。

\subsection{问题提出}

\textbf{问题一：}若每天的电价与小区负载相同，且已知光伏发电功率一天的预测数据，
在微网提供的电能不低于小区负载、储能设备 0:00 与 24:00 储电量相同的条件下，
建立每天 0:00 制定当天计划购电策略的数学模型，并给出指定时段的购电量、
全天购电量与购电费用，以及储能设备的充放电量与 0:00、24:00 的储电量。

\textbf{问题二：}若每天的电价相同，而小区负载与光伏发电功率随时间变化，
当微网提供的电能低于小区负载时，需按交易时刻电价的 5 倍向外网紧急购电。
要求建立每天 0:00 制定计划购电策略的数学模型，并给出指定日期的购电量、
储能设备充放电量与紧急购电情况。

\textbf{问题三：}若每天 0:00、6:00、12:00、18:00 可获得未来 24 小时整点的光伏发电功率预报，
计划购电量与调整购电量之间的偏差分别按交易时刻电价的 50\%（少购）与 150\%（多购）结算，
购电总费用还包括紧急购电费用。要求建立当天计划与调整购电策略的数学模型，
分析是否需要引入其他时刻的预报制定调整策略，并给出指定日期的计算结果。

\textbf{问题四：}实际上外网电价也是实时波动的。要求针对波动电价，
重新建立问题二与问题三条件下的微网购电策略模型，并给出相应的计算结果。
